\exercisesection{}

\begin{enumerate}
\def\labelenumi{\arabic{enumi})}
\tightlist
\item
  Let A be a compact subset of \(\mathbf{C}\) which is equal to the closure of its interior V. We set \(\mathcal{C}(\mathrm{A}) = \mathcal{C}(\mathrm{A};\mathbf{C})\) . Let \(B(A)\) be the vector subspace of \(\mathcal{C}(\mathrm{A})\) consisting of functions \(f\in \mathcal{C}(\mathrm{A})\) whose restriction to V is holomorphic;
\end{enumerate}

it is a Banach subalgebra of \(\mathcal{C}(\mathrm{A})\) . Let \(u \in \mathcal{L}(\mathrm{B}(\mathrm{A}))\) be the linear map such that \(u(f)(z) = zf(z)\) for all \(z \in \mathrm{A}\) and every \(f \in \mathrm{B}(\mathrm{A})\) .

\begin{enumerate}
\def\labelenumi{\alph{enumi})}
\item
  The spectrum of \(u\) is equal to A.
\item
  For every \(z \in V\) the endomorphism u - z of B(A) is a Fredholm endomorphism of index 1.
\item
  For every \(z \in A - V\) the endomorphism u - z of B(A) is not a Fredholm endomorphism.
\item
  The stable spectrum of u is the boundary of A, the sensitive spectrum of u is empty, and the essential spectrum of u is A.
\item
  We assume that the complement of A in \(\mathbf{C}\) is connected. For every \(h \in \mathcal{L}^{c}(\mathrm{B}(\mathrm{A}))\) , the essential spectrum of \(u + h\) is equal to A.
\end{enumerate}

\begin{enumerate}
\def\labelenumi{\arabic{enumi})}
\setcounter{enumi}{1}
\tightlist
\item
  We keep the notation from the previous exercise. We assume that A contains the unit circle \(U\) in \(\mathbf{C}\) and that \(0 \notin A\) . We denote by \(\varphi_{A}\) the characteristic function of A.
\end{enumerate}

\begin{enumerate}
\def\labelenumi{\alph{enumi})}
\tightlist
\item
  The map
\end{enumerate}

\[
\ell \colon f \mapsto \frac {1}{2 i \pi} \int_ {\mathbf {U}} f (z) d z
\]

is a continuous linear form on B(A).

\begin{enumerate}
\def\labelenumi{\alph{enumi})}
\setcounter{enumi}{1}
\tightlist
\item
  Let \(h\) be the finite-rank endomorphism of B(A) such that \(h(f) = \ell(f)\varphi_{\mathrm{A}}\) . The essential spectrum of \(u - h\) is distinct from that of \(u\) . (Show that it contains 0.)
\end{enumerate}

\begin{enumerate}
\def\labelenumi{\arabic{enumi})}
\setcounter{enumi}{2}
\tightlist
\item
  Let E be a complex Hilbert space. Let U be a connected open subset of \(\mathbf{C}\). Let f be a holomorphic function from U into the space \(\mathcal{L}^{c}(\mathrm{E})\) . Suppose that there exists \(z \in U\) such that 1 belongs to the resolvent set of \(f(z)\) .
\end{enumerate}

\begin{enumerate}
\def\labelenumi{\alph{enumi})}
\item
  The set D of \(z \in U\) such that \(1 \in \operatorname{Sp}(f(z))\) is a discrete subset of U.
\item
  The map from U-D to \(\mathcal{L}(\mathrm{E})\) defined by \(z\mapsto \mathrm{R}(f(z),1)\) is holomorphic.
\end{enumerate}

We fix a point \(z_0 \in \mathrm{D}\) .

\begin{enumerate}
\def\labelenumi{\alph{enumi})}
\setcounter{enumi}{2}
\item
  There exists an integer \(k \geqslant 1\) such that the map from U-D into \(\mathcal{L}(\mathrm{E})\) defined by \(z \mapsto (z - z_{0})^{k}\mathrm{R}(f(z), 1)\) extends to a holomorphic map on \((\mathrm{U}-\mathrm{D}) \cup \{z_{0}\}\) .
\item
  There exists a unique sequence \((u_{n})_{n\in\mathbf{N}}\) of endomorphisms of E such that \(u_{n}=0\) for n sufficiently large and a unique holomorphic function S from \((\mathrm{U}-\mathrm{D})\cup\{z_{0}\}\) into \(\mathcal{L}(\mathrm{E})\) such that for all \(z\in U-D\) , we have
\end{enumerate}

\[
\mathrm{R} (f (z), 1) = \sum_ {j = - k} ^ {- 1} (z - z _ {0}) ^ {j} u _ {j} + \mathrm{S} (z).
\]

\begin{enumerate}
\def\labelenumi{\alph{enumi})}
\setcounter{enumi}{4}
\tightlist
\item
  The endomorphism \(u_{-1}\) is the spectral projector of \(f(z_0)\) relative to the open and closed subset \(\{1\}\) of the spectrum of \(f(z_0)\) .
\end{enumerate}

\begin{enumerate}
\def\labelenumi{\arabic{enumi})}
\setcounter{enumi}{3}
\tightlist
\item
  Let E be a Banach space of countable type. Let \(u \in \mathcal{L}(\mathrm{E})\) . Assume that the sequence \((\| u^k \|)_{k \in \mathbf{N}}\) is bounded.
\end{enumerate}

For every \(\lambda\in\mathrm{Sp}(u)\cap\mathbf{U}\) , let \(x_{\lambda}\in\mathrm{E}\) be an eigenvector of \(u\) of norm 1 corresponding to \(\lambda\) .

\begin{enumerate}
\def\labelenumi{\alph{enumi})}
\tightlist
\item
  For all \(\lambda_1\) , \(\lambda_2\) in \(\mathrm{Sp}(u) \cap \mathbf{U}\) and every integer \(k \in \mathbf{N}\) , we have
\end{enumerate}

\[
\left| \lambda_ {1} ^ {k} - \lambda_ {2} ^ {k} \right| \leqslant \left(1 + \| u ^ {k} \|\right) \| x _ {\lambda_ {1}} - x _ {\lambda_ {2}} \|.
\]

\begin{enumerate}
\def\labelenumi{\alph{enumi})}
\setcounter{enumi}{1}
\item
  There exists a real number c \textgreater{} 0 such that \(\|x_{\lambda_{1}} - x_{\lambda_{2}}\| \geqslant c\) for all \(\lambda_{1} \neq \lambda_{2}\) in \(\operatorname{Sp}(u) \cap \mathbf{U}\) .
\item
  The set \(\operatorname{Sp}(u) \cap \mathbf{U}\) is countable (“Jamison’s theorem”).
\item
  Jamison's theorem does not hold in general if E is not of countable type.
\end{enumerate}

\begin{enumerate}
\def\labelenumi{\arabic{enumi})}
\setcounter{enumi}{4}
\tightlist
\item
  The goal of this exercise is to give another proof of Lomonosov's theorem (cor. 2 of III, p. 13) in the case where the space E is a Banach space.
\end{enumerate}

Let E be a complex Banach space of dimension at least 2. Let h be a nonzero compact endomorphism of E and let A be the commutant of h in \(\mathcal{L}(\mathrm{E})\) . a) One of the following conditions is satisfied:

\begin{enumerate}
\def\labelenumi{(\roman{enumi})}
\item
  There exists \(y \in E\) such that Ay is a nonzero subspace that is not dense in E;
\item
  For every closed bounded subset C of A with nonempty interior, there exists a finite subset I of A such that \(h(\mathrm{C})\) is contained in the union of the sets \(u^{-1}(\mathrm{C})\) for \(u \in I\) .
\end{enumerate}

\begin{enumerate}
\def\labelenumi{\alph{enumi})}
\setcounter{enumi}{1}
\tightlist
\item
  Assume that condition (ii) of the previous question is satisfied. Let C be a closed bounded subset of E with nonempty interior. For every integer \(m \geqslant 1\) there exists a finite sequence \((u_{1}, \ldots, u_{m})\) in A such that
\end{enumerate}

\[
u _ {1} \circ \dots \circ u _ {m} \circ h ^ {m} (x _ {0}) \in \mathrm{C}.
\]

(Play table tennis.)

\begin{enumerate}
\def\labelenumi{\alph{enumi})}
\setcounter{enumi}{2}
\item
  Under the same hypothesis, show that the spectral radius of h is not zero (in the contrary case, consider an element \(x_{0} \in E\) such that \(\|h(x_{0})\| > \|h\|\) and the closed ball C with center \(x_{0}\) and radius \(\|h\|\) ; noting that \(\|( \lambda h)^{n}\|\) tends to 0 for every \(\lambda \in R\) , show that the conclusion of the preceding question would imply that 0 belongs to the closure of C).
\item
  For every endomorphism \(u\) of E that commutes with a nonzero compact endomorphism of E, there exists a nonzero closed subspace \(F \subset E\) different from E such that \(u(F) \subset F\) .
\end{enumerate}

\begin{enumerate}
\def\labelenumi{\arabic{enumi})}
\setcounter{enumi}{5}
\tightlist
\item
  Let A be a unital Z-algebra. Assume that A is a free Z-module. Let B be a subset of A that is a basis of A as a Z-module. We say that the pair \((A, B)\) is a natural algebra if the structure constants of A with respect to B (A, III, p. 10) belong to N and if the coefficients of \(1_{A}\) in the basis B belong to N. Moreover, we say that (A, B) is unital if \(1_{A}\) is an element of B.
\end{enumerate}

Let (A, B) be a natural algebra. Let \(B_{0}\) be the set of \(b \in B\) such that the coefficient of b in \(1_{A}\) is nonzero. Denote by \(\tau_{A}\) the unique homomorphism of Z-modules from A to Z such that, for every b in B, one has \(\tau_{\mathrm{A}}(b) = 1\) if \(b \in B_{0}\) , and \(\tau_{\mathrm{A}}(b) = 0\) otherwise.

\begin{enumerate}
\def\labelenumi{\alph{enumi})}
\item
  One has \(\tau_{\mathrm{A}}(xy) = \tau_{\mathrm{A}}(yx)\) for all \((x,y)\in \mathrm{A}^2\) .
\item
  For all \(b \neq b'\) in \(B_{0}\) , we have \(b^2 = b\) and \(bb' = 0\) .
\end{enumerate}

\begin{enumerate}
\def\labelenumi{\arabic{enumi})}
\setcounter{enumi}{6}
\tightlist
\item
  Let (A, B) be a natural algebra.
\end{enumerate}

\begin{enumerate}
\def\labelenumi{\alph{enumi})}
\tightlist
\item
  There exists at most one involution \(\iota\) of \(\mathcal{B}\) such that the unique \(\mathbf{Z}\)-module homomorphism from A to A sending \(b\) to \(\iota(b)\) is an anti-automorphism of A satisfying the condition \(\tau_{\mathrm{A}}(bb') = 1\) if \(b' = \iota(b)\) and \(\tau_{\mathrm{A}}(bb') = 0\) otherwise.
\end{enumerate}

When such an involution \(\iota\) exists, we say that (A, \(\mathcal{B}\) ) is a natural basal algebra.

We now assume that (A, B) is a natural basal algebra.

\begin{enumerate}
\def\labelenumi{\alph{enumi})}
\setcounter{enumi}{1}
\tightlist
\item
  We then have \(\iota(b) = b\) for all \(b \in B_{0}\) as well as the formula
\end{enumerate}

\[
1 _ {\mathrm{A}} = \sum_ {b \in \mathcal {B} _ {0}} b.
\]

\begin{enumerate}
\def\labelenumi{\alph{enumi})}
\setcounter{enumi}{2}
\tightlist
\item
  The structure constants \(\gamma_{bb'}^{b''}\) of A satisfy \(\gamma_{bb'}^{\iota(b'')} = \tau_{\mathrm{A}}(bb'b'')\) for \((b, b', b'')\) in \(\mathcal{B}^{3}\).
\end{enumerate}

\begin{enumerate}
\def\labelenumi{\arabic{enumi})}
\setcounter{enumi}{7}
\tightlist
\item
  Let (A, B) be a natural basal algebra such that A has finite rank.
\end{enumerate}

\begin{enumerate}
\def\labelenumi{\alph{enumi})}
\tightlist
\item
  For every \(x \in A\) , the element
\end{enumerate}

\[
\sum_ {b \in \mathcal {B}} b x \iota (b)
\]

is central.

In the rest of this exercise, we further assume that (A, \(\mathcal{B}\) ) is unital.

\begin{enumerate}
\def\labelenumi{\alph{enumi})}
\setcounter{enumi}{1}
\item
  For all \(b\) and \(b'\) in \(\mathcal{B}\) , there exist \(b_1\) and \(b_2\) in \(\mathcal{B}\) such that the coefficient of \(b'\) in \(bb_1\) and in \(b_2b\) is nonzero.
\item
  For every \(x \in B\) , the matrix representing, in the basis B, the map \(y \mapsto xy\) from A to A is a real matrix all of whose coefficients are positive. We denote by pf the unique homomorphism from A to R such that \(\operatorname{pf}(b)\) is the spectral radius of this matrix when \(b \in B\) .
\end{enumerate}

\(d)\) For every \(b \in \mathcal{B}\) , the real number \(\alpha = \mathrm{pf}(b)\) is an algebraic integer such that \(\alpha \geqslant 1\) ; for every conjugate \(\alpha' \in \mathbf{C}\) of \(\alpha\) , we have \(|\alpha'| \leqslant \alpha\) .

\begin{enumerate}
\def\labelenumi{\alph{enumi})}
\setcounter{enumi}{4}
\tightlist
\item
  Let \(a_0 \in \mathrm{A}\) be the sum of the elements of \(\mathcal{B}\) . All coefficients of the matrix representing the map \(f: y \mapsto ya_0\) in the basis \(\mathcal{B}\) are strictly positive. There exists an element \(r \in \mathbf{C} \otimes \mathrm{A}\) all of whose coefficients in the basis \((1 \otimes b)_{b \in \mathcal{B}}\) are strictly positive such that \(f_{\mathbf{C}}(r) = \varrho(f_{\mathbf{C}})r\) .
\end{enumerate}

\(f)\) For every \(x \in \mathrm{A}\) , there exists a unique complex number \(\delta(x)\) such that \(xr = \delta(x)r\) in \(\mathbf{C} \otimes \mathrm{A}\) . The map \(\delta: \mathrm{A} \to \mathbf{C}\) is a character of \(\mathrm{A}\) .

\begin{enumerate}
\def\labelenumi{\alph{enumi})}
\setcounter{enumi}{6}
\tightlist
\item
  One has \(\delta(x) = \mathrm{pf}(x)\) for all \(x \in \mathrm{A}\) . In particular, the map \(x \mapsto \mathrm{pf}(x)\) is a character of A.
\end{enumerate}

\(h)\) The map pf is the unique character of A such that \(\mathrm{pf}(b) \geqslant 0\) for all \(b \in \mathcal{B}\) ; moreover, one has \(\mathrm{pf}(b) > 0\) for all \(b \in \mathcal{B}\) .

\begin{enumerate}
\def\labelenumi{\roman{enumi})}
\tightlist
\item
  For every \(y \in \mathrm{A}\) , we have \(ry = \delta(y)r\) in \(\mathbf{C} \otimes \mathbf{A}\) .
\end{enumerate}

\begin{enumerate}
\def\labelenumi{\alph{enumi})}
\setcounter{enumi}{9}
\tightlist
\item
  The element
\end{enumerate}

\[
r _ {0} = \sum_ {b \in \mathcal {B}} \mathrm{pf} (b) b
\]

possesses the properties required in question i).

\begin{enumerate}
\def\labelenumi{\arabic{enumi})}
\setcounter{enumi}{8}
\tightlist
\item
  \begin{enumerate}
  \def\labelenumii{\alph{enumii})}
  \tightlist
  \item
    Let \(n\) be an integer and \(M\) a square matrix with coefficients in \(\mathbf{N}\) . One assumes that \(\varrho(M^t M) = \varrho(M)^2\) and that \(|\varrho(M)| < 2\) . There exists an integer \(m \geqslant 2\) such that \(\varrho(M) = 2\cos(\pi/m)\) . (Use Kronecker's theorem, cf. A, V, p. 155, exercise 17.)
  \end{enumerate}
\end{enumerate}

\begin{enumerate}
\def\labelenumi{\alph{enumi})}
\setcounter{enumi}{1}
\item
  For every integer \(m \geqslant 2\) , there exists a square matrix M with coefficients in N such that \(\varrho(M) = 2 \cos(\pi/m)\) .
\item
  Let (A, B) be a unital natural basal algebra of finite rank. Let \(b \in B\) . If \(\operatorname{pf}(b) < 2\) , then there exists an integer \(m \geqslant 3\) such that \(\operatorname{pf}(b) = 2 \cos(\pi/m)\) .
\end{enumerate}

Exercises 10 to 24 provide another proof of Corollary 5 of III, p. 100, as well as certain refinements. The following notions will be used there.

Let I be a finite set and let n be its cardinality.

We denote by \(1_{\mathrm{I}}\) the real identity matrix of type (I, I).

- A square matrix of type (I, I) with real coefficients is said to be positive if all its coefficients are \(\geqslant 0\) and strictly positive if all its coefficients are \(>0\) .

- A square matrix \((a_{ij})_{i,j\in \mathrm{I}}\) of type (I,I) with real coefficients is said to be reducible if there exists a partition of I into two nonempty subsets I' and I'\,' such that \(a_{ij} = 0\) for all \((i,j)\in \mathrm{I}'\times \mathrm{I}''\) ; it is said to be irreducible if it is not reducible.

- One calls a coordinate subspace of \(\mathbf{R}^{\mathrm{I}}\) any vector subspace of \(\mathbf{R}^{\mathrm{I}}\) generated by a subset of the vectors of the canonical basis. There are \(2^{n}\) coordinate subspaces; for each integer \(k\in\{0,\ldots,n\}\) , those among them of dimension \(k\) number \(\binom{n}{k}\) .

A real matrix of type (I, I) is therefore irreducible if and only if the only coordinate subspaces it stabilizes are \(\{0\}\) and \(R^{I}\) .

For all \(x\) and \(y\) in \(\mathbf{R}^{\mathrm{I}}\) we will write \(x \preccurlyeq y\) (and \(y \succcurlyeq x\) ) if \(y - x\) is in the cone of elements of \(\mathbf{R}^{\mathrm{I}}\) with positive coordinates, that is, if \(x_i \leqslant y_i\) for all \(i \in \mathrm{I}\) . We will write \(x \prec y\) (and \(y \succ x\) ) if all coordinates of \(y - x\) are strictly positive.

\begin{enumerate}
\def\labelenumi{\arabic{enumi})}
\setcounter{enumi}{9}
\item
  Let \(A\) be a real matrix of type (I, I), positive and irreducible. For every integer \(n\) such that \(n \geqslant \operatorname{Card}(\mathrm{I})\) , the matrix \((1_{\mathrm{I}} + A)^{n-1}\) is strictly positive. (Let \(y \in \mathbf{R}_{+}^{n}\) and \(z = (1_{\mathrm{I}} + A)y\) ; if \(y \neq 0\), prove that the set of \(i \in \mathrm{I}\) such that \(z_{i} = 0\) is strictly contained in \(\{i \in \mathrm{I} \mid y_{i} = 0\}\) .)
\item
  Let \(A\) be a real matrix of type (I, I), positive and irreducible. We denote by \(\psi \in \mathbf{R}[\mathrm{X}]\) the minimal polynomial of \(A\) and we set \(m = \deg(\psi)\) .
\end{enumerate}

For every integer \(q\), set \(A^q = (a_{ij}^{(q)})_{i,j \in \mathrm{I}}\) .

For every \((i,j)\in\mathrm{I}^{2}\) , there exists an integer q such that \(1\leqslant q\leqslant m\) and \(a_{ij}^{(q)}>0\) . (Let r be the remainder of the Euclidean division of \((1+\mathrm{X})^{n-1}\) by \(\psi\) ; consider the matrix \(r(A)\) and observe that it is strictly positive.)

\begin{enumerate}
\def\labelenumi{\arabic{enumi})}
\setcounter{enumi}{11}
\tightlist
\item
  Let A be a real matrix of type (I, I), irreducible and positive. For every \(x \in R^{I}\) non-zero, we set
\end{enumerate}

\[
r_{x} = \inf_{\substack{i\in \mathrm{I}\\ x_{i}\neq 0}}\frac{(Ax)_{i}}{x_{i}},
\]

where \((Ax)_i = \sum_{k\in I}a_{ik}x_k\) is the coordinate of \(Ax\) with index \(i\) .

\begin{enumerate}
\def\labelenumi{\alph{enumi})}
\item
  If \(x \succcurlyeq 0\) and \(x \neq 0\) , then \(r_x\) is the largest real number \(\varrho\) such that \(\varrho x \preccurlyeq Ax\) .
\item
  Suppose that \(x \succcurlyeq 0\) and set \(y = (1_{\mathrm{I}} + A)x\) . We then have \(r_x \leqslant r_y\) .
\item
  The restriction to \((\mathbf{R}_{+}^{*})^{\mathrm{I}}\) of the function \(x \mapsto r_{x}\) is continuous.
\end{enumerate}

\(d)\) Deduce that there exists \(z \in \mathbf{R}^{\mathrm{I}}\) such that \(z \succcurlyeq 0\) and such that \(r_z = \sup_{x \succcurlyeq 0} r_x\) (Use compact sets \(\mathrm{M} = \{x \in \mathbf{R}_{+}^{\mathrm{I}} \mid \sum_{i \in \mathrm{I}} x_i = 1\}\) as well as \((1_{\mathrm{I}} + A)^{n-1}(\mathrm{M})\) , then exercise 10.)

We set \(r = r_{z}\). An element z of \(R^{I}\) such that \(z \succcurlyeq 0\) and \(r_{z} = r\) will be said to be extremal.

\begin{enumerate}
\def\labelenumi{\alph{enumi})}
\setcounter{enumi}{4}
\item
  We have r \textgreater{} 0.
\item
  If z is an extremal vector, then Az = rz. (If \(u = Az - rz \succcurlyeq 0\) were nonzero, then \((1_{\mathrm{I}} + A)^{n-1}u\) would be \(\succ 0\) based on Exercise 10 and, for \(x = (1_{\mathrm{I}} + A)^{n-1}z\) , one would have \(Ax \succ rx\) .)
\item
  Every extremal element is \(\succ 0\) .
\end{enumerate}

For every vector \(y\) in \(\mathbf{C}^{\mathrm{I}}\) , we denote by \(y^{+}\) the vector \((|y_{i}|)_{i\in \mathrm{I}}\) ; we therefore have \(y^{+}\in \mathbf{R}^{\mathrm{I}}\) and \(y^{+}\succcurlyeq 0\) .

\begin{enumerate}
\def\labelenumi{\alph{enumi})}
\setcounter{enumi}{7}
\item
  Let \(\alpha \in \mathbf{C}\) be an eigenvalue of \(A\) , viewed as a complex matrix, and let \(y\) in \(\mathbf{C}^{\mathrm{I}}\) be a corresponding eigenvector. We have \(|\alpha|y^{+} \preccurlyeq Ay^{+}\) and \(|\alpha| \leqslant r\) .
\item
  Deduce that if \(\alpha = r\) all coordinates of y are nonzero, and then that the eigenspace (in \(C^{I}\) ) of A associated with the eigenvalue r has dimension 1.
\end{enumerate}

We denote by \(\Delta(\mathrm{X})\) the determinant of \(\mathrm{X} \cdot 1_{\mathrm{I}} - A\) and by \(B(\mathrm{X})\) the transpose of the matrix of cofactors of \(\mathrm{X} \cdot 1_{\mathrm{I}} - A\) . We write \(B(\mathrm{X}) = (B_{ij}(\mathrm{X}))_{i,j \in \mathrm{I}}\) . We therefore have \(B(\mathrm{X})(\mathrm{X} \cdot 1_{\mathrm{I}} - A) = (\mathrm{X} \cdot 1_{\mathrm{I}} - A)B(\mathrm{X}) = \Delta(\mathrm{X}) \cdot 1_{\mathrm{E}}\) (A, III, p. 99, prop. 12).

\begin{enumerate}
\def\labelenumi{\alph{enumi})}
\setcounter{enumi}{9}
\tightlist
\item
  For every \(i \in I\) , the column with index i of \(B(r)\) is nonzero. (If this is not the case, construct an eigenvector of A associated with the eigenvalue r whose i-th coefficient is zero, in contradiction with part (f)).
\end{enumerate}

\(k\) ) All the coefficients of \(B(r)\) have the same sign. (First show, using the relation \((r1_{\mathrm{I}} - A)B(r) = 0\) that for all \(i\in \mathrm{I}\) , all the coefficients of the column with index \(i\) of \(B(r)\) are nonzero and of the same sign, and then reason with the transposes.)

\(l)\) One has \(\Delta'(r) > 0\) and \(B_{ij}(r) > 0\) for all \((i,j) \in \mathrm{I}^2\) (Use the formula \(\Delta'(X) = \sum_i B_{ii}(X)\) and observe that \(\Delta'\) does not change sign on the interval \([r, +\infty]\) .)

\begin{enumerate}
\def\labelenumi{\alph{enumi})}
\setcounter{enumi}{12}
\tightlist
\item
  Conclude that r is a simple root of the characteristic polynomial of A.
\end{enumerate}

\begin{enumerate}
\def\labelenumi{\arabic{enumi})}
\setcounter{enumi}{12}
\tightlist
\item
  Let A (resp. C) be a real (resp. complex) square matrix of type (I, I). Assume that A is positive and irreducible, and that \(C^{+} \preccurlyeq A\) , where, similarly to Exercise 12, we denote by \(C^{+}\) the matrix whose coefficients are the absolute values of the coefficients of C, and we write \(C^{+} \preccurlyeq A\) to indicate that the coefficients of \(A - C^{+}\) are positive.
\end{enumerate}

We resume the notation of the preceding exercise.

\begin{enumerate}
\def\labelenumi{\alph{enumi})}
\item
  Let \(\gamma \in \mathbf{C}\) be an eigenvalue of \(C\) and let \(y\) in \(\mathbf{C}^{\mathrm{I}}\) be an associated eigenvector. We have \(|\gamma|y^{+} \preccurlyeq C^{+}y^{+} \preccurlyeq Ay^{+}\) and \(|\gamma| \leqslant r_{y^{+}} \leqslant r\) .
\item
  We further assume that \(|\gamma| = r\) . Then \(y^{+}\) is an extremal vector for A (cf. question f) of Exercise 12), and we have the relations \(Ay^{+} = C^{+}y^{+} = ry^{+}\) and \(A = C^{+}\) .
\end{enumerate}

Let \(u\) (resp. \(u_j\) , for \(j \in I\) ) be the unique complex number such that \(\gamma = |\gamma|u\) (resp. \(y_j = |y_j|u_j\) , for \(j \in I\) ). Let us denote by \(D\) the square complex matrix of type (I, I) whose diagonal coefficients are \((u_j)_{j \in I}\) and whose other coefficients are zero, so that \(y = Dy^+\). Let \(F = u^{-1}D^{-1}CD\) .

\begin{enumerate}
\def\labelenumi{\alph{enumi})}
\setcounter{enumi}{2}
\item
  One has \(Fy^{+} = ry^{+}\) and \(F^{+}y^{+} = Fy^{+}\) (Observe that \(F^{+} = C^{+}\) .)
\item
  One has \(F = F^{+}\) and \(C = uDAD^{-1}\) .
\end{enumerate}

\begin{enumerate}
\def\labelenumi{\arabic{enumi})}
\setcounter{enumi}{13}
\tightlist
\item
  Let A be a real square matrix of type (I, I), positive and irreducible. We resume the notation of Exercise 12.
\end{enumerate}

Let \((\lambda_{h})_{h\in\mathrm{H}}\) be the family of complex eigenvalues of A which have modulus r, repeated with their multiplicities; we set \(v_{h}=\lambda_{h}/r\) for all \(h\in H\) .

Let \(h \in \mathrm{H}\) ; by Exercise 13 applied to \(C = A\) and \(\gamma = \lambda_h\) , there exists a matrix \(D_{(h)}\) such that \(A = v_h D_{(h)} AD_{(h)}^{-1}\) and \(D_{(h)}^+ = 1_{\mathrm{I}}\) .

\begin{enumerate}
\def\labelenumi{\alph{enumi})}
\item
  Let \(z \succ 0\) be an extremal vector of A. For \(h \in H\) , the vector \(y_{(h)} = D_{(h)}z\) is an eigenvector of A for the eigenvalue \(\lambda_{h}\) .
\item
  Show that for \(h \in H\) , the eigenvalue \(\lambda_{h}\) of A is simple. Deduce that the matrices \(D_{(h)}\) and the vectors \(y_{(h)}\) are uniquely determined up to a scalar factor.
\item
  Show that, for all h and k in \(H^{2}\) , we have \(A = u_{h}u_{k}^{-1}D_{(h)}D_{(k)}^{-1}AD_{(k)}D_{(h)}^{-1}\) and hence that the vector \(D_{(h)}D_{(k)}^{-1}z\) is an eigenvector of A associated with the eigenvalue \(u_{h}u_{k}^{-1}r\) .
\item
  Deduce that \(\{u_h\}_{h\in \mathrm{H}}\) is a subgroup of \(\mathbf{C}^*\) and that there exists a family \((\mu_h)_{h\in \mathrm{H}}\) of nonzero complex numbers such that \(\{\mu_h^{-1}D_{(h)}\}_{h\in \mathrm{H}}\) is a subgroup of the group of invertible matrices of type (I,I) with complex coefficients, and that these subgroups are cyclic of order Card(H). (Choose an index \(i\in \mathrm{I}\) and normalize \(\mathrm{D}_{(h)}\) so that its coefficient with row and column index \(i\) is equal to 1.)
\end{enumerate}

In the rest of the exercise, we assume that H = Z/dZ and that there exists \(\omega \in C^{*}\) such that \(\omega^{d} = 1\) and such that \(u_{h} = \omega^{h}\) for \(h \in H\) . Setting \(D = D_{(1)}\) , we then have \(D^{d} = 1_{I}\) and \(D_{(h)} = D^{h}\) for all \(h \in H\) .

\begin{enumerate}
\def\labelenumi{\alph{enumi})}
\setcounter{enumi}{4}
\tightlist
\item
  Deduce that \(A = \omega DAD^{-1}\) .
\end{enumerate}

We consider the partition \((\mathrm{I}_{h})_{h\in\mathrm{H}}\) of I, so that for \(i\in I\) , the coefficient with indices \((i,i)\) of D is equal to \(\omega^{h}\) if and only if \(i\in I_{h}\) . For every h, we set \(n_{h}=\mathrm{Card}(I_{h})\) . The matrix D is written “in blocks”

\[
D = \left( \begin{array}{c c c c} 1 _ {n _ {0}} & & & \\ & \omega 1 _ {n _ {1}} & & \\ & & \ddots & \\ & & & \omega^ {d - 1} 1 _ {n _ {d - 1}} \end{array} \right)
\]

and we consider the corresponding decomposition of \(A\) :

\[
A = \left( \begin{array}{c c c c} A _ {0, 0} & A _ {0, 1} & \dots & A _ {0, d - 1} \\ A _ {1, 0} & A _ {1, 1} & \dots & A _ {1, d - 1} \\ \vdots & \vdots & \ddots & \vdots \\ A _ {d - 1, 0} & A _ {d - 1, 1} & \dots & A _ {d - 1, d - 1} \end{array} \right),
\]

so that for \(h, k \in \mathrm{H}\) , \(A_{h,k}\) is a matrix of type \((\mathrm{I}_h, \mathrm{I}_k)\) .

\(f)\) For all \((h,k)\in\mathrm{H}\) , we have \(\omega A_{h,k}=\omega^{k-h}A_{h,k}\) . Deduce that \(\mathrm{A}_{h,k}=0\) if \(k-h\neq1\) .

\begin{enumerate}
\def\labelenumi{\arabic{enumi})}
\setcounter{enumi}{14}
\tightlist
\item
  We resume the hypotheses and notation of Exercise 12. Let \(\Delta_1(\mathrm{X})\) be the greatest common divisor of the coefficients of \(B(\mathrm{X})\) ; the polynomial \(\Delta_1(\mathrm{X})\) is assumed to be monic.
\end{enumerate}

\begin{enumerate}
\def\labelenumi{\alph{enumi})}
\item
  The polynomial \(\Delta_1(X)\) divides \(\Delta(X)\) and \(\Delta_1(r) \neq 0\) .
\item
  Deduce that the real roots of \(\Delta_{1}(X)\) are strictly smaller than r. Conclude that \(\Delta_{1}(r)>0\) .
\item
  Set \(C(\mathrm{X}) = \frac{1}{\Delta_1(\mathrm{X})} B(\mathrm{X})\) . This is a matrix with coefficients in \(\mathbf{R}[\mathrm{X}]\) . Show that \(C(r) \succ 0\) .
\end{enumerate}

\begin{enumerate}
\def\labelenumi{\arabic{enumi})}
\setcounter{enumi}{15}
\tightlist
\item
  Let A be a real square matrix of type (I, I), positive and irreducible. For all \(j \in I\) set \(s_{j} = \sum_{k \in I} a_{jk}\) ,
\end{enumerate}

\[
s = \inf _ {j \in \mathrm{I}} s _ {j}, \quad \text {and} \quad \mathrm{S} = \sup _ {j \in \mathrm{I}} s _ {j}.
\]

We also adopt the notation \(r\) and \(B(\mathbf{X})\) of Exercise 12.

\begin{enumerate}
\def\labelenumi{\alph{enumi})}
\item
  One has \(\sum_{j\in \mathrm{I}}(r - s_j)B_{kj}(r) = 0\)
\item
  Deduce that \(s \leqslant r \leqslant S\) and that the equality case in either of these two inequalities implies that s = S.
\end{enumerate}

\begin{enumerate}
\def\labelenumi{\arabic{enumi})}
\setcounter{enumi}{16}
\item
  Let \(A\) be a real square matrix of type (I, I), positive and irreducible. Let \(y \in \mathbf{R}^{\mathrm{I}}\) be an eigenvector of \(A\) associated with an eigenvalue \(\alpha\) such that \(y \succcurlyeq 0\) . We then have \(\alpha = r\) (with the notations of Exercise 12) and \(y \succ 0\) . (Consider an extremal vector \(u\) for \(^{t}A\) and compute \(\langle y|^{t}Au\rangle = \langle Ay|u\rangle\) .)
\item
  Let A be a real square matrix of type (I, I), positive and irreducible. Let \(\varrho\) be the spectral radius of A. For every \(x \in R^{I}\) such that \(y \succcurlyeq 0\) and \(x \neq 0\) we set
\end{enumerate}

\[
\varrho^ {x} = \sup _ {i \in \mathrm{I}} \frac {(A x) _ {i}}{x _ {i}},
\]

(with the convention \(\varrho^x = +\infty\) if there exists an index \(i\) such that \(x_i = 0\) and \((Ax)_i \neq 0\) ). Let

\[
\varrho^{\prime} = \inf_{\substack{x\succ 0\\ x\neq 0}}\varrho^{x}.
\]

Following an approach similar to that of Exercise 12, show that \(\varrho'\) is an eigenvalue of \(A\) and that \(\varrho' = \varrho\) .

\begin{enumerate}
\def\labelenumi{\arabic{enumi})}
\setcounter{enumi}{18}
\tightlist
\item
  Let A be a real square matrix of type (I, I), positive and irreducible; let \(\varrho\) be its spectral radius. Let \(z \in R^{I}\) such that \(z \succcurlyeq 0\) and \(z \neq 0\) .
\end{enumerate}

If \(\varrho z \preccurlyeq Az\) (resp. \(\varrho z \succcurlyeq Az\) ), then \(\varrho z = Az\) and \(z \succ 0\) .

\begin{enumerate}
\def\labelenumi{\arabic{enumi})}
\setcounter{enumi}{19}
\tightlist
\item
  Let A be a real square matrix of type (I, I), positive but not necessarily irreducible; let \(\varrho\) be the spectral radius of A. We keep the notations \(B(\mathrm{X})\) and \(C(\mathrm{X})\) from the preceding exercises.
\end{enumerate}

\begin{enumerate}
\def\labelenumi{\alph{enumi})}
\item
  The real number \(\varrho\) is an eigenvalue of \(A\) and there exists at least one eigenvector \(\succcurlyeq 0\) associated with \(\varrho\) .
\item
  For every real number \(\lambda \geqslant r\) , the matrices \(B(\lambda)\) and \(B'(\lambda)\) are positive. (For the second matrix, one may proceed by induction on the cardinality of I.)
\item
  The matrix \(C(\lambda)\) is positive for all \(\lambda \geqslant r\) .
\item
  If A is irreducible, then for every real number \(\lambda > r\) , the matrices \(B(\lambda)\) , \(C(\lambda)\) and \((\lambda1_{\mathrm{I}} - A)^{-1}\) are strictly positive.
\item
  Let J be a subset of I different from I. Let \(\varrho_{\mathrm{J}}\) be the spectral radius of the submatrix \((a_{i,j})_{i,j\in \mathrm{J}}\) . We have \(\varrho_{\mathrm{J}}\leqslant \varrho\) and this inequality is strict if \(A\) is irreducible. Conversely, if \(A\) is reducible, there exists a subset \(\mathrm{J}\subset \mathrm{I}\) different from I, such that \(\varrho_{\mathrm{J}} = \varrho\) .
\item
  Let J be a subset of I different from I and such that the spectral radius of the submatrix \((a_{i,j})_{i,j\in\mathrm{J}}\) is equal to \(\varrho\) . For every subset K such that \(J\subset K\subset I\) , the spectral radius of the matrix \((a_{i,j})_{i,j\in\mathrm{K}}\) is equal to \(\varrho\) . Deduce that \(A\succ0\) if and only if all the diagonal coefficients of \(B(r)\) are strictly positive.
\end{enumerate}

\(g)\) For every real number \(\lambda >\varrho\) and for every subset J of I such that \(\mathrm{J}\neq \mathrm{I}\) the determinant of the matrix \((\lambda \delta_{i,j} - a_{i,j})_{i,j\in \mathrm{J}}\) is strictly positive.

\begin{enumerate}
\def\labelenumi{\alph{enumi})}
\setcounter{enumi}{7}
\tightlist
\item
  Let \(\lambda\) be a real number; we assume that the determinant of \((\lambda \delta_{i,j} - a_{i,j})_{i,j \in \mathbf{J}}\) is strictly positive for every subset \(\mathbf{J} \subset \mathbf{I}\) different from \(\mathbf{I}\) . We then have \(\lambda > \varrho\) .
\end{enumerate}

In the following questions, we assume that \(\mathrm{I}\) is the interval \([1,n]\) in \(\mathbf{N}\). i) Let \((g_{i,j})_{1\leqslant i,j\leqslant n}\) be a real matrix such that \(g_{i,j}\leqslant0\) for all \(i\neq j\) and such that the determinants of the matrices \((g_{i,j})_{1\leqslant i,j\leqslant m}\) are strictly positive for every integer \(m\) such that \(1\leqslant m\leqslant n\) . Then we have \(\det(g_{i,j})_{i,j\in\mathrm{J}}>0\) for every subset J of I.

\begin{enumerate}
\def\labelenumi{\alph{enumi})}
\setcounter{enumi}{9}
\item
  Let \(\lambda\) be a real number; we suppose that for every integer \(m\) such that \(1 \leqslant m \leqslant n\) , the determinant \(\det(\lambda \delta_{i,j} - a_{i,j})_{1 \leqslant i,j \leqslant m}\) is strictly positive. Then \(\lambda > \varrho\) .
\item
  Let \(C = (c_{i,j})_{1 \leqslant i,j \leqslant n}\) be a real matrix such that \(c_{i,j} \geqslant 0\) for all \(i \neq j\) . If \((-1)^{k} \det(c_{i,j})_{1 \leqslant i,j \leqslant m}\) is strictly positive for every integer m such that \(1 \leqslant m \leqslant n\) , then the real parts of the eigenvalues of C (viewed as a complex matrix) are negative.
\end{enumerate}

\begin{enumerate}
\def\labelenumi{\arabic{enumi})}
\setcounter{enumi}{20}
\tightlist
\item
  Let A be a real square matrix of type (I, I), positive; let \(\varrho\) be its spectral radius.
\end{enumerate}

\begin{enumerate}
\def\labelenumi{\alph{enumi})}
\tightlist
\item
  Prove that there exist integers p and q such that \(0 \leqslant q \leqslant p\) and a partition \((\mathrm{I}_{s})_{1 \leqslant s \leqslant p}\) of I into nonempty parts such that the corresponding block decomposition \(A = (A_{s,t})_{1 \leqslant s, t \leqslant p}\) has the following property:
\end{enumerate}

- The matrices \(A_{s,s}\) are irreducible;

- We have \(A_{s,t} = 0\) if \(t > s\) ;

- We have \(A_{s,t} = 0\) if \(1 \leqslant t < s \leqslant q\) ;

- For every integer \(s\) such that \(q \leqslant s \leqslant p\) , there exists \(t\) such that \(A_{s,t} \neq 0\) .

\begin{enumerate}
\def\labelenumi{\alph{enumi})}
\setcounter{enumi}{1}
\item
  Determine all possibilities for such a partition. (In particular, prove that the integers q and p are uniquely determined, as are the sets \(\{I_{1},\ldots,I_{g}\}\) and the sets \(\{I_{q+1},\ldots,I_{p}\}\) .)
\item
  There exists an eigenvector \(z \succ 0\) of A associated with the eigenvalue \(\varrho\) if and only if, for every integer \(s \in [1, p]\) the spectral radius of \(A_{s,s}\) is equal to \(\varrho\) if \(s \leqslant q\) and is strictly less than \(\varrho\) otherwise.
\item
  It is necessary and sufficient for both \(A\) and \({}^t A\) to admit a strictly positive extremal vector that there exists a partition as above such that \(A_{s,t} = 0\) for \(1 \leqslant t < s\) and \(q < s \leqslant p\) and such that the matrices \(A_{s,s}\) for \(s \leqslant q\) have spectral radius equal to \(\varrho\) .
\item
  If A and \(^{t}A\) both admit a strictly positive extremal vector, and if \(\varrho\) has spectral multiplicity equal to 1, then A is irreducible.
\end{enumerate}

\begin{enumerate}
\def\labelenumi{\arabic{enumi})}
\setcounter{enumi}{21}
\tightlist
\item
  Let A be a real square matrix of type (I, I), positive and irreducible; let \(\varrho\) be its spectral radius. Let \(d \geqslant 1\) be the integer determined by Exercise 14. We shall say that A is primitive if d = 1. Let \(\Delta(\mathrm{X})\) be the polynomial \(\det(\mathrm{X} \cdot 1_{\mathrm{I}} - A)\) .
\end{enumerate}

\begin{enumerate}
\def\labelenumi{\alph{enumi})}
\item
  We write \(\Delta(\mathrm{X}) = \mathrm{X}^{n} + a_{1}\mathrm{X}^{n_{1}} + \cdots + a_{t}\mathrm{X}^{n_{t}}\) , where \(a_{1}, \ldots, a_{t}\) are nonzero real numbers and \(n_{1}, \ldots, n_{t}\) natural integers such that \(n > n_{1} > \cdots > n_{t}\) . Prove that d is the greatest common divisor of \(n - n_{1}, n_{1} - n_{2}, \ldots, n_{t-1} - n_{t}\) .
\item
  There exists a nonzero integer m and a polynomial g such that \(\Delta(\mathrm{X}) = g(\mathrm{X}^{d})\mathrm{X}^{m}\) .
\item
  If there exists an integer p \textgreater{} 0 such that \(A^{p}\) is strictly positive, then A is primitive.
\item
  Suppose that \(A\) is primitive. Let us introduce the matrices \(C(\mathbf{X})\) of Exercise 15, and let us set, using the notation of that exercise, \(\mathsf{X}(\mathbf{X}) = \Delta (\mathbf{X}) / \Delta_1(\mathbf{X})\) . The limit of the sequence \((A^k /\varrho^k)_{k\in \mathbf{N}}\) is equal to \(C(\varrho) / \mathsf{X}'(\varrho)\) (One may use a decomposition of \(A\) into Jordan blocks.) Using the results of Exercise 15, deduce that there exists an integer \(p > 0\) such that \(A^p\) is strictly positive.
\item
  Suppose there exists an integer q \textgreater{} 0 such that \(A^{q}\) is reducible. Let p be the greatest common divisor of q and d. Prove that there exists a partition \((\mathrm{I}_{s})_{1 \leqslant s \leqslant p}\) of I into nonempty parts for which \(A^{q}\) is block diagonal, with irreducible diagonal blocks and spectral radius \(\varrho^{q}\) .
\item
  If A is primitive, then \(A^{p}\) is irreducible for every p \textgreater{} 0.
\end{enumerate}

\(g)\) There exists a partition \((\mathrm{I}_s)_{1 \leqslant s \leqslant d}\) of I into nonempty parts for which \(A^d\) is block diagonal, with primitive diagonal blocks and spectral radius equal to \(\varrho^d\) .

\begin{enumerate}
\def\labelenumi{\arabic{enumi})}
\setcounter{enumi}{22}
\tightlist
\item
  Let P be a real square matrix of type (I, I); we say that P is stochastic if it is nonnegative and if the vector \((1, \ldots, 1) \in \mathbf{R}^{\mathrm{I}}\) is an eigenvector of P associated with the eigenvalue 1.
\end{enumerate}

Let A be a real square matrix of type (I, I). Prove that A is positive and admits a strictly positive extremal vector if and only if there exists a stochastic matrix P, a real number r \textgreater{} 0 and a diagonal matrix Z with strictly positive diagonal entries such that \(A = rZPZ^{-1}\) . In this case the spectral radius of \(A\) is equal to \(r\) and an extremal vector of \(A\) is \(Z(1, \ldots, 1)\) .

\begin{enumerate}
\def\labelenumi{\arabic{enumi})}
\setcounter{enumi}{23}
\tightlist
\item
  \begin{enumerate}
  \def\labelenumii{\alph{enumii})}
  \tightlist
  \item
    Let P be a real square stochastic matrix of order n. The real number 1 is a simple root of the minimal polynomial of P. (One may use Exercise 20.) Moreover, every root of modulus 1 of the minimal polynomial of P is simple.
  \end{enumerate}
\end{enumerate}

\begin{enumerate}
\def\labelenumi{\alph{enumi})}
\setcounter{enumi}{1}
\tightlist
\item
  Let A be a positive real matrix admitting an extremal vector \(z \succ 0\) and of spectral radius \(\varrho\) . All roots of modulus \(\varrho\) of the minimal polynomial of A are simple.
\end{enumerate}

In Exercises 25 to 30, we consider a real Fréchet space E and a closed convex cone C in E, assumed to have nonempty interior. The order relation on E induced by C is denoted \(\preccurlyeq_{\mathrm{C}}\) . Thus, for \(x\) , \(y\) in E, we have \(x \succcurlyeq_{\mathrm{C}} y\) if and only if \(x - y \in \mathrm{C}\) . We will write \(x \succ_{\mathrm{C}} y\) for the relation \(x - y \in \mathring{\mathrm{C}}\) .

We consider a continuous map u: E → E satisfying the following conditions:

\begin{enumerate}
\def\labelenumi{(\roman{enumi})}
\item
  The map u is positively homogeneous: we have \(u(\lambda x) = \lambda u(x)\) for every real number \(\lambda \geqslant 0\) ;
\item
  The map u is increasing with respect to C, that is, \(x \succsim_{C} y\) implies \(u(x) \succsim_{C} u(y)\) ;
\item
  For every \(x \in \mathrm{C} - \{0\}\) , we have \(u(x) \in \mathring{\mathrm{C}}\) .
\end{enumerate}

The goal of these exercises, taken from the course of P.-L. Lions at the Collège de France, 2020/2021, is to prove, under certain conditions, that there exists an element \(x \in \mathrm{C} - \{0\}\) and \(\lambda > 0\) such that \(x = \lambda u(x)\) ; this provides in particular another proof of the Krein-Rutman theorem (III, p. 94, th. 4 and III, p. 97, prop. 8).

For \(f \in \mathring{\mathrm{C}}\) and \(\lambda \in \mathbf{R}_{+}\) , one calls an upper solution (resp. lower solution) of the equation \(y = \lambda u(y) + f\) an element \(y \in \mathrm{C}\) such that \(y \succcurlyeq_{\mathrm{C}} \lambda u(y) + f\) (resp. such that \(y \preccurlyeq_{\mathrm{C}} \lambda u(y) + f\) ). We denote by \(U_{f}\) the set of \(\lambda \in \mathbf{R}_{+}\) such that the set of upper solutions of the equation \(y = \lambda u(y) + f\) is nonempty.

\begin{enumerate}
\def\labelenumi{\arabic{enumi})}
\setcounter{enumi}{24}
\tightlist
\item
  \begin{enumerate}
  \def\labelenumii{\alph{enumii})}
  \tightlist
  \item
    Prove that for every \(x \in C\) and every \(y \in \mathring{C}\) , there exists a real number \(\lambda \geqslant 0\) such that \(x \preccurlyeq_{C} \lambda y\) .
  \end{enumerate}
\end{enumerate}

\begin{enumerate}
\def\labelenumi{\alph{enumi})}
\setcounter{enumi}{1}
\tightlist
\item
  The set \(U_{f}\) is an interval that contains 0.
\end{enumerate}

We introduce the following conditions on the map u.

\begin{enumerate}
\def\labelenumi{(\roman{enumi})}
\setcounter{enumi}{3}
\item
  For every \(\lambda \in \mathbf{R}_{+}\) , every \(f \in \mathring{\mathrm{C}}\) , every sequence \((x_{n})\) of elements of C that satisfies \(x_{n+1} = \lambda u(x_{n}) + f\) for all \(n\) and that is either increasing and bounded above, or decreasing, is convergent;
\item
  The image under u of every bounded subset of E is relatively compact.
\item
  For every map \(\lambda \mapsto \varepsilon_{\lambda}\) from \(U_f\) into C such that \(\lim_{\lambda \to \lambda_m} \varepsilon_{\lambda} = 0\) , the set
\end{enumerate}

\[
\{z \in \mathrm{C} \mid \exists \lambda \in \mathrm{U}_{f}\text{ such that } z = \lambda u(z) + \varepsilon_{\lambda} \}
\]

is relatively compact in E.

\begin{enumerate}
\def\labelenumi{\arabic{enumi})}
\setcounter{enumi}{25}
\tightlist
\item
  We resume the preceding notations and hypotheses and we suppose in addition that condition (iv) is satisfied.
\end{enumerate}

\begin{enumerate}
\def\labelenumi{\alph{enumi})}
\item
  Let \(\lambda \in \mathrm{U}_f\) . Let us define a sequence \((x_n)_{n \in \mathbf{N}}\) by \(x_0 = 0\) and \(x_{n+1} = \lambda u(x_n) + f\) for all \(n \in \mathbf{N}\) . It is increasing and bounded above; we have \(x_n \in \mathring{\mathrm{C}}\) for \(n \geqslant 1\) .
\item
  By condition (iv), the sequence \((x_{n})\) converges; let \(x(\lambda)\) denote its limit. Prove that we have \(x(\lambda) \in \mathring{\mathrm{C}}\) and \(x(\lambda) = \lambda u(x(\lambda)) + f\) .
\item
  Let \(\lambda \in \mathrm{U}_f\) and let \(y \in \mathrm{C}\) such that \(y \succcurlyeq_{\mathrm{C}} \lambda u(y) + f\) ; the sequence defined by \(x_0 = y\) , \(x_{n+1} = \lambda u(x_n) + f\) is decreasing, with values in C, and its limit \(x\) satisfies the equation \(x = \lambda u(x) + f\) .
\end{enumerate}

\(d)\) The upper bound of \(\mathrm{U}_{f}\) in \(\mathbf{R}_{+}\) does not belong to \(\mathrm{U}_{f}\) .

\begin{enumerate}
\def\labelenumi{\alph{enumi})}
\setcounter{enumi}{4}
\item
  One has \(U_{g} = U_{f}\) for all \(g \in \mathring{C}\) (For every \(\lambda \in U_{f}\) , let \(x \in C\) be such that \(x \succsim_{C} \lambda u(x) + f\) ; then there exists M \textgreater{} 0 such that y = Mx satisfies \(y \succsim_{C} \lambda u(y) + g\) .) f) Let \(\lambda \leqslant \mu\) be real numbers. Let \(x\) be a lower solution of the equation \(y = \lambda u(y) + f\) and let \(x'\) be an upper solution of the equation \(y = \mu u(y) + f\) . We then have \(x \preccurlyeq_{C} x'\) . (Show that the set of \(\varepsilon \in [0,1]\) such that \(\varepsilon x \preccurlyeq_{C} x'\) coincides with {[}0,1{]}.)
\item
  If C is salient, then for every \(\lambda \in \mathrm{U}_f\) , there exists a unique element \(x \in \mathring{\mathrm{C}}\) such that \(x = \lambda u(x) + f\) .
\end{enumerate}

\(h\) ) Let \(\mu\in\mathbf{R}_{+}\) such that \(f\preccurlyeq_{\mathrm{C}}\mu u(f)\) . We then have \(\mu\notin\mathrm{U}_{f}\) . (The sequence \((x_{n})_{n\in\mathbf{N}}\) defined by \(x_{0}=0\) and \(x_{n+1}=\mu u(x_{n})+f\) satisfies \(x_{n+1}-x_{n}\succcurlyeq_{\mathrm{C}}f\) for all \(n\in\mathbf{N}\) .) In particular, the interval \(\mathrm{U}_{f}\) is bounded. Henceforth we denote by \(\lambda_{\mathrm{m}}\) the supremum of \(\mathrm{U}_{f}\) .

\begin{enumerate}
\def\labelenumi{\roman{enumi})}
\tightlist
\item
  Let \(z \in \mathrm{C} - \{0\}\) and \(\mu \in \mathbf{R}_+\) be such that \(\mu u(z) \succsim_{\mathrm{C}} z\) . We then have \(\lambda_{\mathrm{m}} \leqslant \mu\) , and hence
\end{enumerate}

\[
\lambda_ {\mathrm{m}} \leqslant \inf _ {z \in \mathrm{C} - \{0 \}} \inf \{\mu \mid z \leqslant \mu u (z) \}.
\]

\begin{enumerate}
\def\labelenumi{\alph{enumi})}
\setcounter{enumi}{9}
\tightlist
\item
  The map \(\lambda \mapsto x(\lambda)\) (where \(x(\lambda)\) is defined in the question \(a\) ) is increasing.
\end{enumerate}

\(k)\) The map \(\lambda \mapsto x(\lambda)\) has no limit when \(\lambda\) tends to \(\lambda_{\mathrm{m}}\) . More precisely, if \((\lambda_n)_n\) is a sequence of \(\mathrm{U}_f\) which tends to \(\lambda_{\mathrm{m}}\) , then the sequence \((x(\lambda_n))_n\) is not convergent.

\begin{enumerate}
\def\labelenumi{\arabic{enumi})}
\setcounter{enumi}{26}
\tightlist
\item
  We assume that hypotheses (i) through (v) are satisfied.
\end{enumerate}

\begin{enumerate}
\def\labelenumi{\alph{enumi})}
\item
  Prove that the set of \(x(\lambda)\) , for \(\lambda \in \mathrm{U}_f\) , is not bounded.
\item
  Let d be a metric on E that defines its topology. For \(\lambda \in R_{+}\) , we define \(\hat{x}(\lambda) = d(0, x(\lambda))^{-1}x(\lambda)\) . Prove that the map \(\hat{x}\) is bounded and that the set of values of \(\hat{x}\) has a cluster value \(x_{m}\) such that
\end{enumerate}

\begin{enumerate}
\def\labelenumi{(\roman{enumi})}
\item
  \(d(0, x_{\mathrm{m}}) = 1,\)
\item
  \(x_{\mathrm{m}} = \lambda_{\mathrm{m}} u(x_{\mathrm{m}}),\)
\item
  \(x_{m} \in \overset{\circ}{C}.\)
\end{enumerate}

\begin{enumerate}
\def\labelenumi{\arabic{enumi})}
\setcounter{enumi}{27}
\tightlist
\item
  We assume in this exercise that conditions (i) through (iv) are satisfied. a) We fix \(x_0 \in \mathring{\mathrm{C}}\) . For \(\lambda \in \mathrm{U}_f\) , let \(\mathrm{A}(\lambda) = \inf \{\mathrm{M} \in \mathbf{R}_+ \mid x(\lambda) \preccurlyeq_{\mathrm{C}} \mathrm{M}x_0\}\) . Let \(\mathrm{C}_0 \in \mathbf{R}_+\) be such that \(f \preccurlyeq_{\mathrm{C}} \mathrm{C}_0 u(x_0)\) . For every \(\lambda\) in \(\mathrm{U}_f\) , the interval \([\lambda, \lambda + \mathrm{C}_0 / \mathrm{A}(\lambda)[\) is contained in \(\mathrm{U}_f\) (cf. exercise 26, question \(d\) )).
\end{enumerate}

\begin{enumerate}
\def\labelenumi{\alph{enumi})}
\setcounter{enumi}{1}
\tightlist
\item
  Hence deduce that the function \(\lambda \mapsto \mathrm{A}(\lambda)\) tends to \(+\infty\) as \(\lambda \in \mathrm{U}_f\) tends to \(\lambda_{\mathrm{m}}\) .
\end{enumerate}

For every \(\lambda\) in \(\mathrm{U}_f\) , we define \(\hat{x} (\lambda) = x(\lambda) / \mathrm{A}(\lambda)\) .

We now assume that condition (vi) is satisfied.

\begin{enumerate}
\def\labelenumi{\alph{enumi})}
\setcounter{enumi}{2}
\tightlist
\item
  Prove that the set of values of the mapping \(\hat{x}\) has a cluster value \(x_{\mathrm{m}}\) such that
\end{enumerate}

\[
0 \preccurlyeq_ {\mathrm{C}} x _ {\mathrm{m}} \preccurlyeq_ {\mathrm{C}} x _ {0} \quad \text {and} \quad x _ {\mathrm{m}} = \lambda_ {\mathrm{m}} u (x _ {\mathrm{m}}).
\]

\begin{enumerate}
\def\labelenumi{\alph{enumi})}
\setcounter{enumi}{3}
\item
  Prove that \(x_{m} \neq 0\) . (If one had \(x_{m} = 0\) , let U be a neighborhood of 0 in E such that \(x_{0} + U \subset C\) ; for every \(\mu \geqslant 0\) , there would exist a sequence \((\lambda_{n})\) in \(U_{f}\) converging to \(\lambda_{m}\) such that \(x_{0} - \mu \hat{x}(\lambda_{n}) \succsim_{C} 0\) , contradicting the definition of \(A(\lambda)\) .)
\item
  Let \(y \in \mathrm{C} - \{0\}\) be such that \(y = \lambda_{\mathrm{m}} u(y)\) . Set \(\theta = \sup \{\gamma \mid \gamma y \preccurlyeq_{\mathrm{C}} x_{\mathrm{m}}\}\) . We have \(\theta \in \mathbf{R}_{+}^{*}\) and \(x_{\mathrm{m}} = \theta y\) . (We have \(\theta y \preccurlyeq_{\mathrm{C}} x_{\mathrm{m}}\) and \(\theta y = \lambda_{\mathrm{m}} u(\theta y)\) ; if \(x_{\mathrm{m}} \neq \theta y\) , then we would have \(x_{\mathrm{m}} - \theta y \succ_{\mathrm{C}} 0\) and there would exist \(\varepsilon > 0\) such that \(x_{\mathrm{m}} - \theta y \succcurlyeq_{\mathrm{C}} \varepsilon y\) , contradicting the definition of \(\theta\) .)
\item
  Let \(y \in C - \{0\}\) and \(\mu \in R\) be such that \(y = \mu u(y)\) . We then have \(\mu = \lambda_{m}\) . (Note that \(0 < \mu\) and \(\mu \geqslant \lambda_{m}\) ; if we had \(\mu > \lambda_{m}\) , let \(\theta = \sup\{\gamma \mid \gamma x_{m} \preccurlyeq_{C} y\}\) ; we would then have \(\theta \in R_{+}^{*}\) , then \(\theta x_{m} \preccurlyeq_{C} y\) , whence \(\theta x_{m} = \lambda_{m} u(\theta x_{m}) \prec_{C} \mu u(\theta x_{m}) \preccurlyeq_{C} \mu u(y) = y\) , contradicting the definition of \(\theta\) .)
\end{enumerate}

\begin{enumerate}
\def\labelenumi{\arabic{enumi})}
\setcounter{enumi}{28}
\tightlist
\item
  We resume the notation of the preceding exercise; in particular, we assume that C has nonempty interior. We assume that u is a linear map and that conditions (i), (ii), and (iii) are satisfied.
\end{enumerate}

\begin{enumerate}
\def\labelenumi{\alph{enumi})}
\tightlist
\item
  If u is a compact endomorphism of E, then conditions (iv) and (v) are satisfied.
\end{enumerate}

We assume that conditions (iv) and (vi) are satisfied. Let \((\lambda_{\mathrm{m}}, x_{\mathrm{m}})\) be given by the preceding exercise.

\begin{enumerate}
\def\labelenumi{\alph{enumi})}
\setcounter{enumi}{1}
\tightlist
\item
  For every \(\lambda \in \mathbf{R}_{+}\) such that \(\lambda < \lambda_{\mathrm{m}}\) , the endomorphism \(1_{\mathrm{E}} - \lambda u\) is invertible. (First show that \(1_{\mathrm{E}} - \lambda u\), induced by passage to subspaces, is a bijection of \(\mathring{\mathrm{C}}\) , cf. questions \(a\) ) and \(f\) ) of the preceding exercise; then deduce that
\end{enumerate}

\(1_{E} - \lambda u\) is surjective, and then that it is injective by writing \(x_{1} - x_{2}\) as an element of the kernel of \(1_{E} - \lambda u\) , with \(x_{i} \in \mathring{C}\) , and verifying that there exists \(z \in E\) such that \(x_{1} + z - \lambda u(x_{1} + z) \in \mathring{C}\) .

\begin{enumerate}
\def\labelenumi{\alph{enumi})}
\setcounter{enumi}{2}
\tightlist
\item
  Let y in E be such that \(y = \lambda_{\mathrm{m}} u(y)\) . We then have \(y \in C\) . (There exists c such that \(cx_{m} + y \succ_{C} 0\) ; conclude using part (e) of the preceding exercise.)
\end{enumerate}

\(d)\) Let \(\mathrm{C}^{\circ}\subset\mathrm{E}^{\prime}\) be the polar of C. For every \(x\in\mathring{\mathrm{C}}\) and for every \(\xi\in\mathrm{C}^{\circ}-\{0\}\) , we have \(\langle x,\xi\rangle>0\) .

\begin{enumerate}
\def\labelenumi{\alph{enumi})}
\setcounter{enumi}{4}
\tightlist
\item
  The polar C° of C is preserved by \(^{t}u\colon E'\to E'\) ; for every \(\lambda\in R_{+}\) such that \(\lambda<\lambda_{m}\) , the endomorphism \(1_{E'}-\lambda^{t}u\) of \(E'\) is an isomorphism.
\end{enumerate}

\begin{enumerate}
\def\labelenumi{\arabic{enumi})}
\setcounter{enumi}{29}
\tightlist
\item
  We resume the hypotheses of the preceding exercise, so that u is linear, compact, and conditions (i) to (v) are assumed to be satisfied. The endomorphism \(^{t}u\) is also compact.
\end{enumerate}

\begin{enumerate}
\def\labelenumi{\alph{enumi})}
\item
  Let \(\eta \in \mathrm{C}^{\circ} - \{0\}\) . For every \(\lambda \in \mathbf{R}_{+}\) such that \(\lambda < \lambda_{\mathrm{m}}\) , the unique solution \(\xi(\lambda)\) of the equation \(\xi - \lambda^{t}u(\xi) = \eta\) belongs to \(\mathrm{C}^{\circ}\) . (If \(x \in \mathrm{C}\) , let \(y \in \mathrm{C}\) be such that \(y - \lambda u(y) = x\) ; we have \(\langle y, \eta \rangle = \langle x, \xi \rangle\) .)
\item
  Let \(x_{0} \in \mathring{C}\) . The map \(\lambda \mapsto \langle x_{0}, \xi(\lambda) \rangle\) tends to infinity when \(\lambda\) tends to \(\lambda_{m}\) . Prove that the set of values \(\xi(\lambda)/\langle x_{0}, \xi(\lambda) \rangle\) has a cluster value \(\xi_{m}\) such that:
\end{enumerate}

\begin{enumerate}
\def\labelenumi{(\roman{enumi})}
\item
  \(\xi_{\mathrm{m}}\in \mathrm{C}^{\circ}\) ;
\item
  \(\langle x_{0},\xi_{m}\rangle=1;\)
\item
  \(\xi_{\mathrm{m}} = \lambda_{\mathrm{m}}^{t}u(\xi_{\mathrm{m}})\) .
\end{enumerate}

\begin{enumerate}
\def\labelenumi{\alph{enumi})}
\setcounter{enumi}{2}
\item
  Let \(f \in E\) such that \(\langle f, \xi_{m} \rangle = 0\) . For \(\lambda \in R_{+}\) such that \(\lambda < \lambda_{m}\) , we denote by \(x(\lambda)\) the unique solution \(x\) in E of the equation \(x - \lambda u(x) = f\) . Show that the set of \(x(\lambda)\) is bounded in E (let d be a metric defining the topology of E; if there existed a sequence \(\lambda_{n} \to \lambda_{m}\) such that \(x(\lambda_{n})\) tends to infinity, set \(\hat{x}_{n} = x(\lambda_{n}) / d(0, x(\lambda_{n}))\) , and show that the sequence \((\hat{x}_{n})\) converges to an element \(\hat{x}\) such that \(d(0, \hat{x}) = 1\) , \(\hat{x} = \lambda u(\hat{x})\) and \(\langle \hat{x}, \xi_{m} \rangle = 0\) , which contradicts the fact that one must have \(\hat{x} \in \mathring{C}\) ).
\item
  There exists a unique element \(x \in E\) such that \(f = x - \lambda_{\mathrm{m}} u(x)\) and \(\langle x, \xi_{\mathrm{m}} \rangle = 0\) . (For existence, show using the preceding question that there exists a sequence \((\lambda_{n})\) converging to \(\lambda_{m}\) such that \((x(\lambda_{n}))\) converges to an element x satisfying the required conditions.)
\item
  The endomorphism \(1_{\mathrm{E}'} - \lambda_{\mathrm{m}}{}^{t}u\) of \(\mathrm{E}'\) induces, by passage to subspaces, a bijection of the polar \(x_{\mathrm{m}}^{\circ}\) of \(x_{\mathrm{m}}\) .
\item
  The eigenspace \(\operatorname{Ker}(1_{\mathrm{E}} - \lambda_{\mathrm{m}}u)\) is of dimension 1 and is equal to the polar of the image of \(1_{E} - \lambda_{m}u\) .
\end{enumerate}

CHAPTER IV

